\chapter{Accès~: contrats à terme et options cotées}\label{ch:m1:getting-access-futures-options}

Deux traders achètent le même contrat E-mini au même prix au même
instant. Le premier verse à la bourse 1.18~\$ pour ce privilège~; le second, 35~cents.
Avant même d'avoir la moindre compétence, le second peut tirer profit de
mouvements trois fois plus petits. La différence tient à une adhésion~: un actif ayant un prix de marché, qui
peut être acheté, ou loué au mois à quelqu'un qui en possède une. Les bourses de contrats à terme
publient ouvertement ces écarts~; ils sont le facteur le plus
déterminant de qui peut, tout simplement, mener une stratégie sur contrats à terme à court horizon. Les bourses d'options
ajoutent une couche supplémentaire, la \omterm{def:m1:getting-access-futures-options:appointment}{désignation comme teneur de marché}, qui
permet d'acquérir des privilèges d'allocation autant que des frais réduits. Ce chapitre évalue les deux, et
le contrat qui précède l'un et l'autre~: celui du \omterm{def:m1:getting-access-equities:brokers}{courtier compensateur}.

\section{Le courtier compensateur}

\begin{definition}[Futures commission merchant et membre compensateur]\label{def:m1:getting-access-futures-options:fcm}
Un \emph{futures commission merchant}\index{futures commission merchant} (FCM)
est, aux États-Unis, une entité qui sollicite ou accepte des ordres pour
contrats à terme, options sur contrats à terme ou swaps, et accepte de l'argent ou d'autres actifs
de ses clients pour les couvrir~; il doit être enregistré et membre de
l'organisme d'autorégulation de la profession.
Un \emph{membre compensateur}\index{membre compensateur} est une société agréée par une chambre de compensation pour compenser
des transactions en son nom propre~; elle garantit à la chambre les transactions de ses clients et
contribue au \omterm{def:m1:clearing-and-settlement:waterfall}{fonds de défaillance} (\cref{ch:m1:clearing-and-settlement}).
\end{definition}

Toute position sur contrats à terme est portée par un \omterm{def:m1:getting-access-futures-options:fcm}{membre compensateur}. Une société de trading qui
ne l'est pas est cliente d'un FCM, lequel prélève au moins la
marge de la chambre de compensation (\cref{ch:m1:margin}) et souvent davantage, détient le
collatéral de la société sur des comptes ségrégués, facture une commission par côté et peut
donner à la société accès aux \omterm{def:m1:getting-access-futures-options:membership}{tarifs membres} de ses \omterm{def:m1:getting-access-futures-options:membership}{adhésions à la bourse}, via ses propres
\omterm{def:m1:getting-access-futures-options:incentive}{programmes d'incitation}. Un \emph{give-up}\index{give-up} fonctionne comme sur les actions~: les
transactions exécutées par l'intermédiaire d'un courtier sont remises au FCM compensateur
au titre d'un accord tripartite.

Le FCM supporte le risque de défaillance de la société entre les \omterm{def:m1:financing:margincall}{appels de marge} et le tarifie~:
par le multiple de la marge de bourse qu'il exige, par des limites intrajournalières sur
la position et la taille des ordres, par le droit de liquidation. Pour une stratégie
à levier, ce sont ces conditions, et non la commission, qui constituent le coût de la relation.
En périodes de tension (\cref{ch:m1:stress-case-studies}), les FCM relèvent leurs
multiples et écartent les clients dont le risque de queue leur déplaît~; l'accès d'une société
aux contrats à terme n'est aussi sûr que l'appétit de son FCM.

\begin{omfigure}
\begin{tikzpicture}[font=\small,
  box/.style={draw=omDef,rounded corners=2pt,minimum width=2.7cm,minimum height=0.95cm,align=center,fill=omDef!4},
  ext/.style={draw=omExo,rounded corners=2pt,minimum width=2.7cm,minimum height=0.95cm,align=center,fill=omExo!6},
  flow/.style={-{Stealth[length=2mm]},thick}]
\node[box] (f) at (0,0) {Société de trading};
\node[box] (fcm) at (5.4,0) {FCM compensateur\\(marge, limites)};
\node[ext] (ch) at (10.8,0) {Chambre de compensation};
\node[ext] (ex) at (10.8,2.0) {Bourse};
\node[box] (eb) at (5.4,2.0) {Courtier exécutant\\ou adhésion propre};
\draw[flow,omDef] (f.north) |- node[pos=0.75,above,font=\footnotesize]{ordres} (eb.west);
\draw[flow,omDef] (eb) -- (ex);
\draw[flow,omThm] (ex) -- node[right,font=\footnotesize]{transactions} (ch);
\draw[flow,omThm] ([yshift=4pt]ch.west) -- node[above,font=\footnotesize]{appels} ([yshift=4pt]fcm.east);
\draw[flow,omThm] ([yshift=4pt]fcm.west) -- node[above,font=\footnotesize]{ses appels} ([yshift=4pt]f.east);
\draw[flow,omExo,dashed] (eb) -- node[right,font=\footnotesize]{give-up} (fcm);
\end{tikzpicture}

\omcaption{Accès à une bourse de contrats à terme. L'adhésion modifie les frais sur le
chemin supérieur~; elle ne modifie pas le chemin inférieur~: sauf si la société devient
\omterm{def:m1:getting-access-futures-options:fcm}{membre compensateur}, un FCM se place entre elle et la chambre de compensation et fixe
sa marge réelle.}\label{fig:m1:getting-access-futures-options:chain}
\end{omfigure}

\section{Adhésion~: acheter, louer ou rester à l'écart}

\begin{definition}[Adhésion à la bourse et tarif membre]\label{def:m1:getting-access-futures-options:membership}
Une \emph{adhésion à la bourse}\index{adhésion à la bourse} d'une bourse de contrats à terme
est un droit cessible, détenu par un particulier ou une société, qui donne accès au
\emph{tarif membre}\index{tarif membre}~: un barème de frais de bourse inférieur à celui
appliqué à tous les autres. Les adhésions sont réparties par groupe de produits, peuvent
s'acheter et se revendre à des prix publiés par la bourse, et se louer
au mois.
\end{definition}

\begin{dated}{2026-09}{Un contrat, trois prix}\label{dat:m1:getting-access-futures-options:rates}
Un comparatif publié par un courtier en contrats à terme indique, pour l'E-mini S\&P 500, des
frais de bourse par côté de 1.18~\$ pour un non-membre, 0.47~\$ pour le locataire
d'une adhésion, 0.35~\$ pour le propriétaire~; pour le \omterm{def:m1:index-and-single-stock-futures-worldwide:point}{contrat micro}, 0.20~\$, 0.07~\$ et
0.04~\$. Les non-membres paient aussi des frais réglementaires de 0.02~\$ par côté. La location
est une procédure en deux étapes~: l'agrément comme membre personne physique de la bourse (des
frais de dossier de 2\,000~\$), puis un contrat de location avec le titulaire d'un siège.
La bourse a annoncé des modifications de ses frais de transaction applicables à compter du
1\up{er}~octobre 2026~; son barème de frais fait foi.
\end{dated}

\begin{omfigure}
\begin{tikzpicture}
\begin{axis}[width=9.5cm,height=4.8cm,ybar stacked,bar width=18pt,
  ylabel={\$ par côté},
  symbolic x coords={non-member,lessee,owner},xtick={non-member,lessee,owner},xticklabels={non-membre,locataire,propriétaire},
  tick label style={font=\small},label style={font=\small},xticklabel style={font=\small\strut},
  ymin=0,ymax=1.7,ymajorgrids,grid style={gray!25},enlarge x limits=0.25,
  legend columns=3,legend style={font=\footnotesize,at={(0.5,-0.2)},anchor=north,draw=none,fill=none,/tikz/every even column/.append style={column sep=8pt}},
  table/col sep=comma]
\addplot[fill=omDef!60,draw=omDef] table[x=status,y=exchange]{figdata/markets-1/30-getting-access-futures-options/per_side.csv};
\addplot[fill=omThm!60,draw=omThm] table[x=status,y=regulatory]{figdata/markets-1/30-getting-access-futures-options/per_side.csv};
\addplot[fill=omExo!50,draw=omExo] table[x=status,y=commission]{figdata/markets-1/30-getting-access-futures-options/per_side.csv};
\legend{frais de bourse,frais réglementaires,commission du courtier}
\end{axis}
\end{tikzpicture}

\omcaption{Ce que coûte un côté d'un E-mini~: 1.45~\$, 0.57~\$ et 0.45~\$. Frais de bourse
et frais réglementaires issus de l'\cref{dat:m1:getting-access-futures-options:rates}~;
les commissions, 25 cents pour les non-membres et 10 pour les membres, sont
illustratives. Données~: le tutoriel.}\label{fig:m1:getting-access-futures-options:perside}
\end{omfigure}

\begin{proposition}[Volume de rentabilité]\label{prop:m1:getting-access-futures-options:breakeven}
Si une voie d'accès à la bourse coûte $c_1$ par côté, avec un coût fixe mensuel $K_1$,
et une autre $c_2 < c_1$ avec $K_2 > K_1$, la seconde est moins chère au-delà de
\[
n^* = \frac{K_2 - K_1}{c_1 - c_2}\ \text{côtés par mois.}
\]
\end{proposition}

Avec une location à 1\,500~\$ par mois (une hypothèse~; les loyers sont cotés au mois et
varient selon la division et la demande), l'économie de 88~cents par côté couvre la
location à partir de 1\,705 côtés par mois, soit environ 40 allers-retours par jour. La propriété immobilise
le prix de l'adhésion~; avec un coût du capital de 4\,000~\$ par mois, elle devient plus avantageuse que la
location à partir de 20\,833 côtés par mois, soit mille par jour
(\cref{fig:m1:getting-access-futures-options:monthly}).

\begin{omfigure}
\begin{tikzpicture}
\begin{axis}[width=11cm,height=5cm,
  xlabel={côtés par mois (milliers)},ylabel={milliers de \$ par mois},
  tick label style={font=\small},label style={font=\small},
  xmin=0,xmax=30,ymin=0,ymax=30,grid=major,grid style={gray!25},
  legend columns=3,legend style={font=\footnotesize,at={(0.5,-0.3)},anchor=north,draw=none,fill=none,/tikz/every even column/.append style={column sep=8pt}},
  table/col sep=comma]
\addplot[omThm,very thick] table[x=sides_k,y=non_member_k]{figdata/markets-1/30-getting-access-futures-options/monthly.csv};
\addplot[omDef,very thick,dashed] table[x=sides_k,y=lessee_k]{figdata/markets-1/30-getting-access-futures-options/monthly.csv};
\addplot[black,very thick,dashdotted] table[x=sides_k,y=owner_k]{figdata/markets-1/30-getting-access-futures-options/monthly.csv};
\legend{non-membre,locataire,propriétaire}
\end{axis}
\end{tikzpicture}

\omcaption{Coût mensuel de la négociation de l'E-mini selon les trois voies. La
location est amortie presque immédiatement~; la propriété ne l'est que pour une société qui négocie
environ mille côtés par jour. Données~: le tutoriel.}\label{fig:m1:getting-access-futures-options:monthly}
\end{omfigure}

Les frais comptent en proportion du tick. Un aller-retour au tarif de non-membre
représente 19~\% d'un tick de l'E-mini et 35~\% d'un tick du micro
(\cref{fig:m1:getting-access-futures-options:ticks})~: une stratégie qui vise à capter un seul tick
conserve les quatre cinquièmes de celui-ci sur le grand contrat et moins de deux
tiers sur le petit, avant toute commission. Le petit contrat est le
plus cher, rapporté à ce qu'il peut rapporter, pour tout le monde~; il n'est bon marché que pour
les membres.

\begin{omfigure}
\begin{tikzpicture}
\begin{axis}[width=10cm,height=4.6cm,ybar,bar width=12pt,
  ylabel={\% d'un tick},
  symbolic x coords={non-member,lessee,owner},xtick={non-member,lessee,owner},xticklabels={non-membre,locataire,propriétaire},
  tick label style={font=\small},label style={font=\small},xticklabel style={font=\small\strut},
  ymin=0,ymax=42,ymajorgrids,grid style={gray!25},enlarge x limits=0.25,
  nodes near coords,nodes near coords style={font=\footnotesize,fill=white,inner sep=1pt,/pgf/number format/.cd,fixed,precision=0},
  legend columns=2,legend style={font=\footnotesize,at={(0.5,-0.2)},anchor=north,draw=none,fill=none,/tikz/every even column/.append style={column sep=8pt}},
  table/col sep=comma]
\addplot[fill=omDef!60,draw=omDef] table[x=status,y=es_pct]{figdata/markets-1/30-getting-access-futures-options/ticks.csv};
\addplot[fill=omThm!60,draw=omThm] table[x=status,y=micro_pct]{figdata/markets-1/30-getting-access-futures-options/ticks.csv};
\legend{E-mini (tick 12.50~\$),micro (tick 1.25~\$)}
\end{axis}
\end{tikzpicture}

\omcaption{Frais de bourse et frais réglementaires d'un aller-retour, en pourcentage
d'un tick. Données~: \cref{dat:m1:getting-access-futures-options:rates}.}\label{fig:m1:getting-access-futures-options:ticks}
\end{omfigure}

\section{Programmes d'incitation et programmes de teneur de marché}

\begin{definition}[Programme d'incitation]\label{def:m1:getting-access-futures-options:incentive}
Un \emph{programme d'incitation}\index{programme d'incitation} (incentive programme) est un barème de remises ou de rabais sur les frais
que propose une bourse de contrats à terme, pour un produit ou un
groupe de produits et pour une durée limitée, aux participants qui atteignent des seuils de volume
ou qui respectent des \omterm{def:m1:options-market-structure:mm}{obligations de cotation}~; l'adhésion n'est souvent pas exigée.
\end{definition}

Les bourses de contrats à terme ne versent pas de rabais aux fournisseurs de liquidité de façon généralisée, comme le font les bourses d'actions~;
elles les ciblent. Les nouveaux produits, les produits en concurrence avec ceux
d'une autre bourse, les échéances éloignées et les heures de nuit bénéficient de programmes de \omterm{def:m1:what-a-trading-firm-does:market-maker}{teneur de marché}
assortis d'obligations (largeur maximale, taille minimale, pourcentage
du temps) et, en contrepartie, des indemnités ou des exonérations de frais~; les sociétés de trading
à fort volume obtiennent des remises échelonnées. Les conditions sont déposées auprès du régulateur, mais
la participation se fait sur demande, les places peuvent être limitées, et le programme
peut prendre fin. Une société dont la stratégie n'est rentable que dans le cadre d'un programme a
une activité assortie d'une échéance.

\section{Désignations comme teneur de marché d'options}

\begin{definition}[Désignation comme teneur de marché]\label{def:m1:getting-access-futures-options:appointment}
Une \emph{désignation comme teneur de marché}\index{désignation comme teneur de marché} (market-maker appointment) est
l'enregistrement d'un membre, sur une bourse d'options, comme \omterm{def:m1:what-a-trading-firm-does:market-maker}{teneur de marché} pour un ensemble de
classes d'options. Elle emporte les \omterm{def:m1:options-market-structure:mm}{obligations de cotation} et les privilèges de la
\cref{def:m1:options-market-structure:mm}, et se paie au moyen de permis de négociation et de
frais de désignation qui croissent avec le nombre et l'activité
des classes.
\end{definition}

Sur une bourse d'options, l'écart de frais entre les \omterm{def:m1:what-a-trading-firm-does:market-maker}{teneurs de marché} et les autres
est plus large que sur les contrats à terme, et il s'accompagne des droits d'allocation décrits au
\cref{ch:m1:options-market-structure}. Le côté des coûts est aussi plus lourd~: des permis
sur chacune d'une douzaine de bourses ou davantage, puisque la même série se négocie sur toutes~;
la bande passante et les ports pour la cotation en masse~; le flux consolidé de
l'\cref{dat:m1:options-market-structure:penny}~; et du capital pour un inventaire d'options
qui ne peut pas être soldé à la clôture. La tenue de marché d'options a une taille minimale efficace
que n'a pas le trading pour compte propre sur contrats à terme, ce qui explique
qu'elle soit assurée par un petit nombre de sociétés.

\section{Licences de données}

Les couches sont celles du \cref{ch:m1:getting-access-equities}~: accès, affichage,
usage hors affichage, redistribution. Les contrats à terme ajoutent une distinction qui surprend les
nouvelles sociétés~: sur certaines bourses, les frais de données en temps réel sont supprimés ou réduits pour
les membres qui négocient pour leur propre compte, et facturés en totalité aux non-membres et à
tout usage qui ressemble à un service rendu à d'autres. Le statut qui réduit les frais de transaction
réduit aussi la facture de données, et le seuil de rentabilité de la
\cref{prop:m1:getting-access-futures-options:breakeven} doit tenir compte des deux.

\section{Tutoriel~: quand une location devient-elle rentable~?}\label{tut:m1:getting-access-futures-options:lease}

\begin{tutorial}
\textbf{Objectif.} Calculer le coût total d'un côté selon trois voies, les
volumes de seuil de rentabilité entre elles, l'effet d'un \omterm{def:m1:getting-access-futures-options:incentive}{programme d'incitation} au volume, et le
coût d'un aller-retour en ticks. \textbf{État final~:} les trois chiffres
de données de ce chapitre.
\begin{tutsteps}
\item \textbf{Un côté.} Frais de bourse et frais réglementaires par statut, plus la
commission du courtier.
\omcode{code/firm/futfees/firm_futfees.py}{10}{29}{Frais par statut, voies et coût d'un mois.}\label{lst:m1:getting-access-futures-options:fees}
\item \textbf{Seuil de rentabilité.}
\omcode{code/firm/futfees/firm_futfees.py}{32}{38}{Nombre mensuel de côtés au-delà duquel la voie au coût fixe le plus élevé l'emporte.}\label{lst:m1:getting-access-futures-options:breakeven}
\item \textbf{Incitations} qui s'appliquent à l'ensemble du mois dès qu'un seuil est
franchi~: les effets de palier du \cref{ch:m1:getting-access-equities} une fois encore.
\end{tutsteps}
\textbf{Pour aller plus loin.} Ajouter la facture de données, avec l'exonération des membres, au
coût fixe de chaque voie. Ajouter ensuite un second produit dans une autre division et
déterminer si une seule adhésion ou deux revient moins cher pour une répartition donnée du volume.
\end{tutorial}

\section{Construction~: le moteur de frais des contrats à terme}\label{bld:m1:getting-access-futures-options:futfees}

\begin{build}
\bfield{Rôle} Le backtest facture chaque exécution simulée à ce que la
société paierait réellement~; chaque trimestre, la direction financière décide
de louer, d'acheter ou d'abandonner des adhésions.
\bfield{Interface} \texttt{ProductFees(product, exchange\_fee, regulatory\_fee)}
par statut~; \texttt{Access(status, broker\_commission, fixed\_monthly)}; \texttt{per\_side}, \texttt{monthly\_cost}, \texttt{breakeven\_sides(product,
cheap, dear)}~; \texttt{IncentiveTier(min\_sides, discount)} et
\texttt{with\_incentive}; \texttt{cost\_in\_ticks(product, access, tick\_value)}.
\bfield{Règles} Les frais sont des données à dates d'effet, indexées par produit et
statut. Les coûts fixes incluent la location, ou le coût du capital engagé pour
une adhésion achetée. Une voie plus chère par côté n'a jamais de seuil de rentabilité~: la
fonction renvoie rien, et non un nombre négatif. Les \omterm{def:m1:futures-contracts-and-exchanges:tick}{valeurs du tick} proviennent de la
référence des contrats (\cref{bld:m1:futures-contracts-and-exchanges:master}).
\bfield{Tests d'acceptation} \texttt{code/firm/futfees/tests/}~: coût par côté
et par mois~; seuil de rentabilité dans les deux sens~; une incitation qui s'applique à
l'ensemble du mois à son seuil~; frais par rapport au tick.
\bfield{Pour aller plus loin} Options sur contrats à terme et options cotées, avec des frais par
capacité (client, professionnel, société, \omterm{def:m1:what-a-trading-firm-does:market-maker}{teneur de marché}) et selon que
l'ordre a ajouté ou retiré de la liquidité.
\end{build}

\begin{omsources}
\item National Futures Association, \emph{Futures commission merchant (FCM)
registration}.
\item Optimus Futures, ``CME membership seat lease'' (comparaison publiée des
frais de bourse des non-membres, des locataires et des propriétaires).
\item CME Group, \emph{Exchange fees for clearing and trading} et
\emph{Membership pricing} (pages de référence~; non lisibles par machine à la date
de rédaction).
\end{omsources}

\section{Exercices}

\begin{exercise}[$\star$]\label{exo:m1:getting-access-futures-options:1}
Avec les frais de la \cref{fig:m1:getting-access-futures-options:perside}, donner
le coût tout compris de 10\,000 côtés en un mois pour un non-membre et pour un
locataire versant 1\,500~\$ pour la location.
\end{exercise}

\begin{exercise}[$\star$]\label{exo:m1:getting-access-futures-options:2}
Calculer les deux volumes de seuil de rentabilité cités dans le texte~: louer
plutôt que rester à l'extérieur, et posséder (à 4\,000~\$ par mois de coût du capital) face à
la location.
\end{exercise}

\begin{exercise}[$\star$]\label{exo:m1:getting-access-futures-options:3}
Exprimer les frais de bourse et réglementaires d'un aller-retour en pourcentage
d'un tick pour l'E-mini (tick 12.50~\$) et le \omterm{def:m1:index-and-single-stock-futures-worldwide:point}{contrat micro} (tick 1.25~\$), pour
chaque statut.
\end{exercise}

\begin{exercise}[$\star\star$]\label{exo:m1:getting-access-futures-options:4}
Une stratégie réalise 400 allers-retours par jour sur l'E-mini et gagne, avant frais,
en moyenne 0.3 tick par aller-retour. Donner son P\&L mensuel (21 jours) en tant
que non-membre, puis en tant que locataire, avec les commissions de la figure.
\end{exercise}

\begin{exercise}[$\star\star$]\label{exo:m1:getting-access-futures-options:5}
Une bourse accorde une remise de 5 cents par côté à partir de 50\,000 côtés par mois et de 10
cents à partir de 200\,000, sur tous les côtés du mois. Un locataire a négocié 46\,000
côtés avec deux jours restants, et il en négocie normalement 1\,500 par jour. Quelle est
la valeur des 1\,000 derniers côtés~?
\end{exercise}

\begin{exercise}[$\star\star$]\label{exo:m1:getting-access-futures-options:6}
Un FCM demande 150~\% de la marge de bourse et laisse deux heures pour couvrir les
\omterm{def:m1:margin:intraday}{appels de marge intrajournaliers}. Un autre en demande 110~\% et dix minutes. Pour une stratégie
à fortes positions intrajournalières et petites positions de nuit, laquelle est la plus chère~? Que
demander d'autre aux deux~?
\end{exercise}

\begin{exercise}[$\star\star\star$]\label{exo:m1:getting-access-futures-options:7}
\emph{Programmation.} Ajouter un coût mensuel de données de 1\,200~\$ pour les non-membres, exonéré pour
les locataires et les propriétaires, aux voies du tutoriel. Recalculer le seuil de
rentabilité entre rester à l'extérieur et louer. Que devient-il si la location
vaut 1\,100~\$~?
\end{exercise}

\begin{exercise}[$\star\star\star$]\label{exo:m1:getting-access-futures-options:8}
\emph{Trouver la faille.} «~Notre stratégie de market making sur un contrat à terme nouvellement coté
rapporte 180\,000~\$ par mois~: 30\,000~\$ proviennent de la fourchette et 150\,000~\$ de
la rémunération de \omterm{def:m1:what-a-trading-firm-does:market-maker}{teneur de marché} versée par la bourse. Nous levons des capitaux pour
multiplier cette activité par dix.~» Quelle question un investisseur devrait-il poser~?
\end{exercise}

\section{Problème~: louer ou payer~?}

\begin{problem}[{Problème du week-end --- première adhésion d'une petite société}]\label{pb:m1:getting-access-futures-options:1}
Une société de deux personnes négocie le E-mini par l'intermédiaire d'un FCM en tant que non-membre~: 600
côtés par jour, 21 jours par mois, aux frais de
la \cref{fig:m1:getting-access-futures-options:perside}. Son bénéfice brut de trading
avant frais s'élève en moyenne à 0.25 tick par côté (\omterm{def:m1:futures-contracts-and-exchanges:tick}{valeur du tick}~: 12.50~\$). Une location est
proposée à 1\,500~\$ par mois~; en tant que locataire, sa commission tombe à 10 cents.

\medskip\noindent\textbf{Partie I --- Aujourd'hui.}
\begin{enumerate}
\item Donner le nombre de côtés par mois et le bénéfice brut par mois.
\item Donner les frais mensuels en tant que non-membre.
\item Donner le bénéfice net, et les frais exprimés en pourcentage du bénéfice brut.
\item Donner le coût d'un aller-retour exprimé en ticks.
\end{enumerate}

\medskip\noindent\textbf{Partie II --- Avec une location.}
\begin{enumerate}[resume]
\item Donner le volume de seuil de rentabilité en côtés par mois et par jour.
\item Donner les frais mensuels en tant que locataire, location comprise.
\item Donner le nouveau bénéfice net et l'amélioration obtenue.
\item La demande d'adhésion coûte 2\,000~\$ une seule fois. Combien de jours d'économies faut-il pour la rembourser~?
\item Quels sont les coûts non financiers liés au fait de devenir membre~?
\end{enumerate}

\medskip\noindent\textbf{Partie III --- Croissance.}
\begin{enumerate}[resume]
\item À partir de quel volume la possession, au coût de capital de 4\,000~\$ par mois, serait-elle
plus rentable que la location~?
\item La société ajoute le \omterm{def:m1:index-and-single-stock-futures-worldwide:point}{contrat micro}, 2\,000 côtés par jour à 0.25 tick brut
(\omterm{def:m1:futures-contracts-and-exchanges:tick}{valeur du tick}~: 1.25~\$). Donner son bénéfice net mensuel en tant que non-membre et en tant que locataire, avec
des commissions de 10 et 5 cents.
\item Que dit la question 11 sur les produits qu'un non-membre peut négocier
à court terme~?
\item Un \omterm{def:m1:getting-access-futures-options:incentive}{programme d'incitation} accorde une remise de 10 cents par côté au-delà de 20\,000 côtés
par mois. Donner sa valeur pour la société à son volume sur le E-mini.
\item Le programme prend fin le trimestre prochain. Comment la société doit-elle traiter sa valeur dans
ses plans~?
\end{enumerate}

\medskip\noindent\textbf{Partie IV --- Jugement.}
\begin{enumerate}[resume]
\item Pourquoi les bourses facturent-elles trois fois plus cher aux non-membres~?
\item Qui détient des adhésions pour les louer, et pourquoi s'agit-il d'une activité à part entière~?
\item Le FCM de la société relève son multiple de marge de 120~\% à 200~\% après un
mois très volatil. Qu'est-ce qui pèse le plus pour la société, cela ou la location~?
\item Qu'est-ce qui pousserait la société à envisager de devenir elle-même \omterm{def:m1:getting-access-futures-options:fcm}{membre compensateur}~?
\item Énoncer le \emph{résultat nommé}~: le volume de seuil de rentabilité de la location en
côtés par mois.
\item En une phrase~: qu'est-ce qu'une adhésion permet d'acheter~?
\end{enumerate}
\end{problem}

\section{Questions d'entretien}

\begin{interviewq}[$\star$ \iqroles{trader, développeur}]\label{iq:m1:getting-access-futures-options:1}
Qu'est-ce qu'un FCM et à quoi sert-il pour une société de trading~?
\end{interviewq}

\begin{interviewq}[$\star$ \iqroles{trader, chercheur}]\label{iq:m1:getting-access-futures-options:2}
Pourquoi les frais des bourses de contrats à terme diffèrent-ils entre membres et non-membres, et
dans quelle mesure~?
\end{interviewq}

\begin{interviewq}[$\star\star$ \iqroles{chercheur, trader}]\label{iq:m1:getting-access-futures-options:3}
Votre backtest d'une stratégie de scalping à un tick sur un \omterm{def:m1:index-and-single-stock-futures-worldwide:point}{contrat micro} est
rentable avant frais. Que vérifiez-vous ensuite~?
\end{interviewq}

\begin{interviewq}[$\star\star$ \iqroles{trader}]\label{iq:m1:getting-access-futures-options:4}
Quand est-il judicieux de louer une adhésion, et quand est-il préférable d'en acheter une~?
\end{interviewq}

\begin{interviewq}[$\star\star$ \iqroles{trader, banque}]\label{iq:m1:getting-access-futures-options:5}
Que négocieriez-vous avec un \omterm{def:m1:getting-access-equities:brokers}{courtier compensateur}, et qu'est-ce qui
n'est pas négociable~?
\end{interviewq}

\begin{interviewq}[$\star\star\star$ \iqroles{trader, chercheur}]\label{iq:m1:getting-access-futures-options:6}
Pourquoi le market making d'options cotées est-il concentré dans quelques sociétés, alors que
le trading pour compte propre sur contrats à terme ne l'est pas~?
\end{interviewq}
